\chapter{Worldlines: Execution as Geodesics}
\label{chapter-6-worldlines-execution-as-geodesics}

\section{What It Means for a Computation to Move}
\label{what-it-means-for-a-computation-to-move}

Up to this point we have lived in possibility. We charted the cone of immediate futures (Kairos), the
laws constraining them (Aios), and the fossilized past (Chronos). We built a static map of a
dynamically evolving universe. We stood outside time, looking at what \emph{could} happen.

But engineers don't get paid for possibility. We ship software that \emph{actually} runs. We care
about the specific, singular trajectory our system takes through the dark sea of Rulial Space. We
care about the one path that executes, not the million paths that didn't. Every other worldline
remains a ghost—an untraversed branch in the multiverse of what might have been.

\begin{quote}
\textbf{Which path does the computation actually walk?}
\end{quote}

That path—the fully realized, irreversible sequence of rewrites from input to output—is called a
\textbf{worldline}. It is not a metaphor borrowed from physics. It is not poetic language draped
over dry machinery. In \COMPUTER{}, it is a rigorously defined mathematical object with
measurable properties. This chapter formalizes it, demystifies it, and shows how every act of
debugging and optimization is nothing but geometry applied to these lines.

When we say a computation "moves," we mean something precise: the universe state changes. One graph
becomes another. Nodes are deleted, edges rewired, new structure appears. The \WARP graph that was
$S_0$ a moment ago is now $S_1$, and the transformation between them happened because some rule
fired, some match was selected, some pattern was consumed and replaced. That step—that singular,
irreversible rewrite—is the atom of computational motion. String enough of these atoms together, and
you get a worldline: the complete trajectory of a running program through the geometry of Rulial Space.

This is fundamentally different from how we typically think about execution. In classical models,
execution is a black box. The CPU does something. Instructions happen. State mutates somewhere in
RAM. We instrument it with logs and breakpoints, but the structure of the path itself remains
invisible. In \COMPUTER{}, the path \emph{is} the structure. Every log entry you've ever read,
every stack trace you've chased, every bug you've fought—they all live exclusively on this line. The
worldline \emph{is} the computation. Everything else is commentary.

\section{What Is a Worldline?}
\label{sec:worldline-definition}

A worldline is a totally ordered sequence of legal DPO rewrites applied under a deterministic
Scheduler. The traditional vocabulary—"running code," "executing instructions," "the program
counter"—is too fuzzy to describe what actually happens. In WARP space, execution is nothing but an
unbroken chain of graph states:

\[
S_0 \to S_1 \to S_2 \to \dots \to S_N.
\]

Each arrow is a rewrite. Each state $S_i$ is a complete snapshot of the universe at tick $i$. There
are no gaps, no hidden transitions, no magic. If you want to know what the program did between $S_5$
and $S_6$, the answer is: exactly one rewrite, chosen by the Scheduler, applied via DPO. If you want
to know what the entire execution did, the answer is: this sequence, in this order, deterministically.

This bears repeating because it is so alien to how we usually think. A worldline does not summarize
execution. It does not approximate it. It \emph{is} execution, with perfect fidelity. Every
intermediate state is materialized. Every choice is recorded. Every branching moment is collapsed
into a single canonical decision. The worldline is the truth.

%----------------------------------------------------------------------
% NEW DIAGRAM: Worldline as Path Through State Space
%----------------------------------------------------------------------
\begin{figure}[h]
\centering
\begin{tikzpicture}[
    state/.style={circle, draw=blue!60, fill=blue!10, minimum size=8mm, drop shadow},
    path/.style={-{Stealth}, thick, blue!80!black},
    ghost/.style={circle, draw=gray!60, fill=gray!15, minimum size=6mm, dashed},
    ghostpath/.style={-{Stealth}, thin, gray!60, dashed},
    >=Stealth
]

% Cloud of possibility (background) - DRAW FIRST so it's behind everything
\draw[gray!30, rounded corners, dashed, fill=gray!5]
    (-0.5,-1.8) rectangle (8.5,2.2);
\node[font=\scriptsize, gray] at (7.5,1.9) {Rulial Space};

% The worldline path (bold, followed)
\node[state] (s0) at (0,0) {$S_0$};
\node[state] (s1) at (2,0.5) {$S_1$};
\node[state] (s2) at (4,0.3) {$S_2$};
\node[state] (s3) at (6,-0.2) {$S_3$};
\node[state] (s4) at (8,0.4) {$S_4$};

\draw[path] (s0) -- (s1);
\draw[path] (s1) -- (s2);
\draw[path] (s2) -- (s3);
\draw[path] (s3) -- (s4);

% Ghost paths (not taken)
\node[ghost] (g1) at (2,-1.2) {};
\node[ghost] (g2) at (4,1.5) {};
\node[ghost] (g3) at (6,1.8) {};

\draw[ghostpath] (s0) -- (g1);
\draw[ghostpath] (s1) -- (g2);
\draw[ghostpath] (s2) -- (g3);

% Labels
\node[font=\footnotesize, below] at (s0.south) {Initial};
\node[font=\footnotesize, above] at (s4.north) {Halted};
\node[font=\footnotesize, gray, below] at (g1.south) {Not taken};
\node[font=\footnotesize, gray, above] at (g2.north) {Not taken};

\end{tikzpicture}
\caption{A worldline is the actual path taken through Rulial Space. Ghost paths represent
possible futures that were never chosen. The Scheduler collapses this cloud of possibility into a
single deterministic sequence.}
\label{fig:worldline-path}
\end{figure}

%----------------------------------------------------------------------
% Math vs Runtime Analogy Table
%----------------------------------------------------------------------
\begin{table}[h]
\centering
\begin{tabularx}{\linewidth}{@{}l l X@{}}
\toprule
\textbf{Concept} & \textbf{Mathematical Object} & \textbf{Runtime Intuition} \\
\midrule
State & WARP $S$ & Heap / object graph at a tick \\
Rule & DPO production $(L,K,R)$ & Instruction / function with pre- and post-conditions \\
Match & Monomorphism $L \hookrightarrow S$ & Concrete pattern match on in-memory objects \\
Scheduler & Selection function $\sigma(S)$ & Interpreter loop / thread scheduler \\
Worldline & Sequence $(S_i)_{i=0}^N$ & Execution trace / commit history of the universe \\
Rulial Distance & Metric $d(S,T)$ & ``How many steps?'' between two states \\
Geodesic & Distance minimizing path & Best possible implementation / optimal hot path \\
Collapse & Choice of one match & Branch decision / control flow edge taken \\
Neighborhood & Local ball in $d$ & Set of easily reachable behaviors / refactor targets \\
\bottomrule
\end{tabularx}
\caption{Mathematical objects in the \COMPUTER{} model and their runtime analogies.}
\label{table:math-vs-runtime}
\end{table}

\begin{center}
\rule{0.5\linewidth}{0.5pt}
\end{center}

\section{Why Worldlines Are Deterministic}
\label{6.2-why-worldlines-are-deterministic}

Classical graph rewriting is nondeterministic. If three rules match, the theory says: "Pick one."
Maybe all three, in parallel universes. Fun for theorists—unusable for engineers. You cannot
reproduce a bug if the system picks different rules each time. You cannot optimize what you cannot
measure. You cannot reason about what refuses to sit still.

\COMPUTER{} rejects nondeterminism outright. We impose a ruthless authoritarian:
\textbf{The Scheduler}. It does not guess. It does not roll dice. It does not consult a random
number generator or check the wall clock. Every time it looks at a state $S$, it computes the same
canonical next step. Run the same input twice, you get the same worldline. Byte-for-byte,
edge-for-edge, tick-for-tick. This is not a nice-to-have property. It is the foundation on which
everything else stands.

To understand \emph{how}, we must expose the mechanism behind the determinism.

\subsection{The Mechanics of Determinism: How the Scheduler Breaks Ties}

For any given state $S$, there may be dozens of valid rule matches. The left-hand side of rule $A$
might match nodes $(n_1, n_2)$. The left-hand side of rule $B$ might match nodes $(n_2, n_3)$. Both
are legal. Both satisfy the gluing conditions. Both could fire. Classical graph rewriting shrugs and
says "either one." \COMPUTER{} does not shrug. It makes a choice, and that choice is always the
same.

The Scheduler collapses this ambiguity using a strict, layered decision stack. First, it enumerates
all legal matches as tuples of node IDs, edge IDs, and port bindings. Because the WARP graph has
stable identities—IDs that do not depend on memory addresses, allocation order, or the phase of the
moon—these tuples form an ordered set. Second, it sorts the matches lexicographically by rule ID,
then by match tuple shape, then by node and edge identity order. Third, overlapping matches are
eliminated via a deterministic leftmost-priority filter: if two matches touch the same node or edge,
the one that appears earlier in the sorted order wins, and the other is discarded. Finally, if
multiple viable matches remain—a rare case, but possible—the Scheduler computes a canonical
embedding hash and selects the match minimizing it. This effectively performs a topological sort over
match space, breaking the final ties with a stable, reproducible function.

%----------------------------------------------------------------------
% TikZ: Scheduler Tie-Breaker Visualization
%----------------------------------------------------------------------

\begin{figure}[h]
\centering
\begin{tikzpicture}[
    node/.style={circle, draw, minimum size=7mm, inner sep=1pt, drop shadow},
    matchA/.style={rounded corners, draw=blue, thick, dashed},
    matchB/.style={rounded corners, draw=orange!80!black, thick, dashed},
    >=Stealth
]

% Nodes (universe state S)
\node[node] (n1) at (0,0) {$n_1$};
\node[node] (n2) at (2.5,0) {$n_2$};
\node[node] (n3) at (5,0) {$n_3$};

% Edges
\draw[-{Stealth}] (n1) -- node[above, font=\scriptsize] {$e_{12}$} (n2);
\draw[-{Stealth}] (n2) -- node[above, font=\scriptsize] {$e_{23}$} (n3);

% Highlight match A: (n1, n2)
\node[matchA, fit=(n1) (n2), inner sep=10pt, label={[blue]above:$m_A$}] (mA) {};

% Highlight match B: (n2, n3)
\node[matchB, fit=(n2) (n3), inner sep=10pt, label={[orange!80!black]below:$m_B$}] (mB) {};

% Annotations
\node[align=left, font=\footnotesize] at (0,-1.5)
    {$\mathsf{ids}(m_A) = (n_1, n_2)$};
\node[align=left, font=\footnotesize] at (5,-1.5)
    {$\mathsf{ids}(m_B) = (n_2, n_3)$};

\node[align=center, font=\footnotesize, anchor=west] at (-0.5,1.6)
    {Same rule ID: $\mathrm{id}(r_A) = \mathrm{id}(r_B)$};

% Arrow showing lexicographic preference
\draw[->, thick] (0.2,-2.0) -- (4.8,-2.0)
    node[midway, below, font=\footnotesize]
    {$\mathsf{ids}(m_A) <_{\text{lex}} \mathsf{ids}(m_B)$, so $m_A$ wins};

\end{tikzpicture}
\caption{Deterministic tie-breaking: two competing matches of the same rule.
Canonical ID ordering selects $m_A$ as the unique next rewrite.}
\label{fig:scheduler-tie-breaker}
\end{figure}

This hierarchy ensures that the transition $S_i \rightarrow S_{i+1}$ is both deterministic and
\textbf{canonical}. Replay the same input, you get the same universe. Byte-for-byte, edge-for-edge.
No race conditions. No timing bugs. No Heisenbugs that vanish when you attach a debugger. The system
has one legal future, and the Scheduler finds it every single time.

This property is the bedrock on which everything else stands. Without it, we have nothing. (For the
formal specification of how this machinery works, see Section~\ref{sec:scheduler-spec} at the end of
this chapter.)

\begin{center}
\rule{0.5\linewidth}{0.5pt}
\end{center}

\section{Worldline Sharpness: Sensitivity to Initial Conditions}
\label{6.3-worldline-sharpness}

Two universes can share the same first hundred ticks yet diverge catastrophically at tick 101. A
single flipped bit, a reordered port, or a one-edge difference can alter which rules match at the
next step. By tick 200, the worldlines may not even resemble each other. One may have halted
successfully. The other may be stuck in an infinite loop, allocating memory until the heat death of
the system.

This sensitivity is called \textbf{Worldline Sharpness}. It is the computational analog of chaos in
dynamical systems: small changes in initial conditions lead to large differences in outcomes. But
unlike chaotic systems, where sensitivity is often treated as a nuisance, in \COMPUTER{} it is a
feature we can measure and reason about. Systems with low curvature resist deviation; errors dampen
out, the path self-corrects, and small mistakes are forgiven by the geometry. Systems with high
curvature amplify deviation; one wrong rewrite can throw the worldline into the abyss, and recovery
becomes impossible.

Every engineer who's debugged a race condition has experienced high-curvature geometry firsthand.
You change one line of code—a single function call, moved up or down by three lines—and the entire
system behaves differently. Not slightly differently. Completely differently. The bug vanishes. Or it
gets worse. Or it moves to a different module entirely. This is not magic. It is geometry. The
worldline bent, and you are now walking a different path through Rulial Space.

%----------------------------------------------------------------------
% NEW DIAGRAM: Worldline Divergence / Sharpness
%----------------------------------------------------------------------
\begin{figure}[h]
\centering
\begin{tikzpicture}[
    state/.style={circle, draw=black!60, fill=black!10, minimum size=6mm, drop shadow},
    path/.style={-{Stealth}, thick},
    >=Stealth
]

% Common initial path
\node[state] (s0) at (0,0) {$S_0$};
\node[state] (s1) at (1.5,0) {$S_1$};
\node[state] (s2) at (3,0) {$S_2$};

\draw[path, blue!80!black] (s0) -- (s1) -- (s2);

% Divergence point label
\node[font=\footnotesize, above] at (s2.north) {Divergence};

% Upper path (worldline A)
\node[state, fill=blue!20] (a3) at (4.5,1.2) {$A_3$};
\node[state, fill=blue!20] (a4) at (6,1.8) {$A_4$};
\node[state, fill=blue!20] (a5) at (7.5,2.2) {$A_5$};

\draw[path, blue!80!black] (s2) -- (a3) -- (a4) -- (a5);
\node[font=\scriptsize, blue!80!black, above right] at (a5.east) {Worldline A};

% Lower path (worldline B)
\node[state, fill=orange!20] (b3) at (4.5,-1.2) {$B_3$};
\node[state, fill=orange!20] (b4) at (6,-1.8) {$B_4$};
\node[state, fill=orange!20] (b5) at (7.5,-2.2) {$B_5$};

\draw[path, orange!80!black] (s2) -- (b3) -- (b4) -- (b5);
\node[font=\scriptsize, orange!80!black, below right] at (b5.east) {Worldline B};

% Annotation
\node[font=\footnotesize, align=center] at (5.5,0)
    {Tiny change at $S_2$\\leads to completely\\different futures};

% Distance measure
\draw[<->, thick, red!70!black] (a5.south) -- (b5.north)
    node[midway, right, font=\scriptsize] {$d(A_5, B_5) \gg 0$};

\end{tikzpicture}
\caption{Worldline sharpness: a small perturbation at $S_2$ causes worldlines A and B to diverge
dramatically. High-sharpness systems amplify tiny differences into catastrophic outcomes.}
\label{fig:worldline-divergence}
\end{figure}

Sharpness is not inherently bad. Sometimes you want sensitivity. A security system should react
violently to unauthorized input. A compiler should reject malformed syntax immediately, not try to
limp along. But when sharpness is accidental—when it emerges from poorly structured rules or tangled
dependencies—it becomes a nightmare. Debugging high-sharpness systems feels like defusing a bomb:
touch anything, and the universe explodes.

\begin{center}
\rule{0.5\linewidth}{0.5pt}
\end{center}

\section{Geodesics: The ``Straight Lines'' of Computation}
\label{6.4-geodesics}

If Rulial Distance measures the minimum number of rewrites required to transform $A$ into $B$, then a
\textbf{computational geodesic} is the worldline that achieves that minimum. It is the straightest
path through curved space. It is the fastest route from input to output. It is the worldline with no
wasted motion.

\begin{quote}
\textbf{The geodesic is the most efficient execution possible.}
\end{quote}

Real software, of course, rarely follows the geodesic. We wander. We allocate memory and free it a
moment later. We iterate empty lists. We check conditions we already know. We call functions that do
nothing. We build data structures we never read. These detours are not abstract inefficiencies—they
are literal curvature in the path. Every unnecessary rewrite bends the worldline away from the
geodesic. Every wasted cycle is a deviation that could have been avoided.

Why does this happen? Sometimes it is deliberate. We add logging, bounds checking, defensive copies.
We trade performance for safety or debuggability. But often, the curvature is accidental. It emerges
from poor design, accumulated cruft, or ignorance of better algorithms. The code does too much
because nobody ever stopped to ask: what is the minimum path from here to there?

%----------------------------------------------------------------------
% NEW DIAGRAM: Geodesic vs Curved Worldline
%----------------------------------------------------------------------
\begin{figure}[h]
\centering
\begin{tikzpicture}[
    state/.style={circle, draw, fill=gray!10, minimum size=5mm, drop shadow},
    geodesic/.style={-{Stealth}, thick, green!60!black, line width=1.5pt},
    curved/.style={-{Stealth}, thick, red!70!black},
    >=Stealth
]

% Start and end states - moved down by 1.5 units
\node[state, fill=green!20] (start) at (0,-2.0) {$S_{\text{in}}$};
\node[state, fill=green!20] (end) at (8,-2.0) {$S_{\text{out}}$};

% Geodesic (straight line)
\draw[geodesic] (start) -- (end)
    node[midway, above, font=\footnotesize, green!60!black] {Geodesic: $d_{\min}$ steps};

% Curved worldline (with detours) - moved down by 1.5 units
\node[state, fill=red!10] (c1) at (2,-0.8) {};
\node[state, fill=red!10] (c2) at (3,-0.2) {};
\node[state, fill=red!10] (c3) at (5,-0.5) {};
\node[state, fill=red!10] (c4) at (6.5,-1.2) {};

\draw[curved] (start) to[out=45, in=180] (c1);
\draw[curved] (c1) to[out=0, in=180] (c2);
\draw[curved] (c2) to[out=0, in=120] (c3);
\draw[curved] (c3) to[out=-60, in=120] (c4);
\draw[curved] (c4) to[out=-60, in=90] (end);

\node[font=\footnotesize, red!70!black, align=center] at (4,0.4)
    {Actual worldline:\\detours, waste, curvature};

% Distance annotation - moved down by 1.5 units
\node[font=\scriptsize, below] at (4,-2.7)
    {Optimization = reducing curvature, approaching the geodesic};

\end{tikzpicture}
\caption{The geodesic (green) is the shortest path through Rulial Space. Real worldlines (red)
wander, accumulating unnecessary rewrites. Optimization is the process of straightening the path.}
\label{fig:geodesic-comparison}
\end{figure}

The geodesic is not always obvious. In a sufficiently complex system, finding it may be computationally
intractable. But we can often find \emph{approximations}—worldlines that are straighter than what we
started with, even if not perfectly straight. This is what profilers and optimizers do: they identify
the parts of the worldline where curvature is highest, and they suggest ways to reduce it. Inline
this function. Memoize that result. Hoist this computation out of the loop. Each of these
transformations is a geometric operation: unbending the path, reducing the deviation from the
geodesic.

Understanding geodesics changes how we think about performance. It is not about "making things
faster" in some vague sense. It is about \emph{aligning the worldline with the geometry of the
space}. Fast code feels elegant because it is geometrically direct. Slow code feels wrong because it
is wandering through high-curvature regions for no good reason.

\subsection{Worldline Inertia: Curvature as Technical Debt}

Curved worldlines possess \textbf{inertia}: once bent, they resist straightening. Every detour adds
structure to the universe. New structure alters which rules can fire in the future. Altered rule
availability reshapes the geometry ahead. What began as a small inefficiency—a single unnecessary
allocation, say—grows into a dependency. Other parts of the code start assuming that allocation
exists. Functions are written to consume it. Invariants are established around it. The longer the
curved worldline persists, the more mass it gathers, and the harder it becomes to straighten.

This is the geometric definition of \textbf{technical debt}. Bad decisions add curvature to the
worldline. Curvature gathers mass as the system evolves around it. Mass bends the local geometry,
making it harder for future rewrites to take the straight path. Bending the geometry increases the
cost of correction. Eventually, the debt becomes so entrenched that paying it off requires rewriting
large portions of the system—not because the code is "bad" in some abstract sense, but because the
\emph{geometry of the worldline} has become irreparably warped.

Refactoring is not merely changing code. It is unbending spacetime. You are reshaping the geometry so
that future worldlines can take straighter paths. You are reducing the mass of accumulated cruft so
that the system can move more freely. This is why good refactors feel like relief: you are literally
straightening out the universe.

\subsection{Rulial Friction: Why Some Paths Hurt More Than Others}

Inertia explains why curved worldlines resist straightening. But curvature alone
is not the whole story. Some parts of the universe are easy to rewrite—light,
loose, low-stakes. You can shuffle them around with minimal consequences. Others are load-bearing
megastructures: dense dependency hubs, core libraries, highly-shared subgraphs. Touch one node, and
fifty others feel the ripple. These regions carry what we call \textbf{Rulial Friction}.

Low-friction regions rewrite cleanly. A small nudge realigns the path. You refactor a helper
function, and nothing else notices. High-friction regions resist every attempt at correction. Even
tiny edits require rewriting large amounts of connected structure. You touch a core abstraction, and
the entire codebase convulses.

A hacky shortcut often works by tunneling straight through a high-friction region and perturbing
nodes that were never meant to move. It "works" in the sense that it produces the desired output, but
it leaves the geometry scarred. Later, when you attempt a proper refactor, the cost is enormous: the
graph itself pushes back. Every change you try to make conflicts with the dependencies that formed
around the hack. You are not fighting the code. You are fighting the geometry.

\begin{quote}
Technical debt is not just curvature. It is curvature through high-friction space.
\end{quote}

This is why some refactors feel effortless and others feel geological. The graph itself
pushes back against changes that violate its load-bearing structure.

%----------------------------------------------------------------------
% TikZ: Rulial Friction Diagram
%----------------------------------------------------------------------
\begin{figure}[h]
\centering
\begin{tikzpicture}[
    node/.style={circle, draw=black, minimum size=6mm, inner sep=1pt, drop shadow},
    lightnode/.style={circle, draw=black!40, fill=black!10, minimum size=6mm, drop shadow},
    heavynode/.style={circle, draw=black!70, fill=gray!40, minimum size=6mm, drop shadow},
    edge/.style={-{Stealth}, thin},
    heavyedge/.style={-{Stealth}, thick},
    dashedarrow/.style={->, dashed, thick, blue},
    >=Stealth
]

% Low friction nodes (left)
\node[lightnode] (l1) at (-2,0.8) {};
\node[lightnode] (l2) at (-3,0) {};
\node[lightnode] (l3) at (-1,0) {};
\node[lightnode] (l4) at (-2,-0.8) {};

\draw[edge] (l1) -- (l2);
\draw[edge] (l1) -- (l3);
\draw[edge] (l2) -- (l4);
\draw[edge] (l3) -- (l4);

\node at (-2,-1.4) {\footnotesize Low Friction Region};

% High friction cluster (right)
\node[heavynode] (h1) at (2,1) {};
\node[heavynode] (h2) at (1.3,0.3) {};
\node[heavynode] (h3) at (2.7,0.3) {};
\node[heavynode] (h4) at (1.5,-0.6) {};
\node[heavynode] (h5) at (2.5,-0.6) {};

\foreach \a/\b in {h1/h2, h1/h3, h2/h3, h2/h4, h3/h5, h4/h5, h2/h5, h3/h4}
    \draw[heavyedge] (\a) -- (\b);

\node at (2,-1.4) {\footnotesize High Friction Region};

% Worldline arrows
\draw[dashedarrow, blue] (-4,0) -- (-1,0);
\node[font=\scriptsize, blue, below] at (-2.5,-0.3) {Easy refactor};

\draw[dashedarrow, red!80!black] (0,0) .. controls (1,1.5) .. (3,1)
    node[pos=0.6, above, font=\scriptsize, yshift=2pt] {Hard refactor (high friction)};

\end{tikzpicture}
\caption{Rulial Friction: rewriting a sparse region is easy (low friction),
but rewrites in dense, load-bearing regions require massive structural change.
Technical debt is not just curvature—it is curvature through high-friction space.}
\label{fig:rulial-friction}
\end{figure}

\begin{mdframed}[linewidth=1pt, linecolor=black!50, innertopmargin=8pt, innerbottommargin=8pt]
\noindent\textbf{FOR THE NERDS™: A Formal Definition of Rulial Friction}

\medskip

Let $S$ be a universe state, and let $m = (r, \mu)$ be a legal match producing successor
state $S' = r(S)$.

We define the \textbf{local friction measure} $\Phi(S, m)$ as:

\[
\Phi(S, m)
= \left| \{ x \in S \;|\; x \notin S' \text{ or } x \text{ has altered adjacency in } S' \} \right|.
\]

That is: the number of nodes and edges whose structural role changes under the rewrite.

Intuitively, if a rewrite touches only a few nodes (low $|\Phi|$), the region has low friction. If a rewrite perturbs a dense subgraph (high $|\Phi|$), the region has high friction.

We extend this to a directional derivative along a worldline $(S_i)$:

\[
\mathcal{F}(S_i)
= \Phi(S_i, \sigma(S_i)),
\]

the friction encountered at tick $i$. Finally, define the \textbf{integrated friction} over a worldline segment:

\[
\mathcal{F}_{a \to b}
= \sum_{i=a}^{b} \mathcal{F}(S_i).
\]

This is the Rulial analog of "work" in physics:
the cumulative structural resistance encountered by the computation.

\begin{quote}
Low-friction universes evolve cleanly.
High-friction universes fight every rewrite.
\end{quote}

\end{mdframed}

\begin{mdframed}[linewidth=1pt, linecolor=black!50, innertopmargin=8pt, innerbottommargin=8pt]
\noindent\textbf{FOR THE NERDS™: Curvature, Friction, and Refactor Difficulty}

\medskip

\textit{(A full treatment of curvature appears in Chapter 9.
This subsection provides only the local, first-order precursor.)}

Friction measures how many elements a rewrite perturbs. To capture how \emph{sensitive}
a region is to those perturbations, we introduce a notion of \textbf{local curvature}.

Let $S$ be a state and $m = (r,\mu)$ a match. Consider a small \emph{perturbation family}
of matches $\{m'\}$ differing slightly from $m$ (e.g.\ by reattaching one edge, toggling a port,
or modifying a local subgraph). Let $S' = r(S)$ and $S'_{m'} = r_{m'}(S)$ be the corresponding
successor states.

Define the \textbf{local curvature} at $(S,m)$ as:

\[
\kappa(S,m)
= \max_{m' \in \mathcal{N}(m)} d\big(S', S'_{m'}\big),
\]

where $\mathcal{N}(m)$ is a small neighborhood of perturbed matches and
$d$ is the Rulial Distance.

If $\kappa$ is small, then nearby choices lead to similar futures—a flat region where mistakes are
cheap. If $\kappa$ is large, then tiny changes lead to wildly different futures—high curvature, where
one wrong move can send the worldline careening into disaster.

We then define the \textbf{Refactor Difficulty} measure:

\[
\Psi(S,m) = \kappa(S,m) \cdot \Phi(S,m),
\]

the product of curvature ($\kappa$) and friction ($\Phi$) at a point. This gives us four regimes.
Low friction and low curvature: easy edits, forgiving futures—the sweet spot of clean code. High
friction but low curvature: heavy but predictable—costly but safe, like refactoring a large,
well-structured library. Low friction but high curvature: light but chaotic—subtle changes with
explosive consequences, the domain of race conditions and timing bugs. Finally, high friction
\emph{and} high curvature: dense structure with explosive sensitivity—the absolute worst case, where
every attempt to fix something breaks three other things.

This last regime is what engineers experience as "touch-anything-and-everything-breaks." It is not a
feeling. It is not bad luck. It is a measurable property of the geometry. You are trying to perform
brain surgery on a live, flailing universe, and the universe is not cooperating.

\begin{quote}
Refactor Difficulty is not a feeling. It is a measurable property of the geometry.
\end{quote}

\end{mdframed}

\begin{center}
\rule{0.5\linewidth}{0.5pt}
\end{center}

\section{Collapse: Choosing One Future}
\label{6.5-collapse}

The Time Cone presents a cloud of potential futures. The worldline requires a single committed one.
That transition is \textbf{Collapse}—the moment the Scheduler prunes the entire possibility tree,
keeps one branch, and discards the rest.

Collapse is not mystical. It is not quantum. It is not pulling some metaphysical trick borrowed from
physics we don't understand. It is mechanical. It is the filtering of all legal matches down to a
single canonical rewrite. Given a state $S$ and a set of matching rules, the Scheduler evaluates
them, orders them, eliminates conflicts, and picks the winner. That winner is the rule that fires.
Everything else is forgotten.

Every conditional branch in imperative code is a local Collapse. The \texttt{if} statement evaluates
a boolean, and one of two paths is taken. The other ceases to exist. Every loop guard is a Collapse:
continue or break. Every function call is a Collapse: which implementation gets invoked? Every
pattern match is a Collapse: which arm of the \texttt{match} expression wins? Collapse is the act
that converts possibility into history. It is the instant where the future becomes the past, and the
present ceases to be an open question.

%----------------------------------------------------------------------
% NEW DIAGRAM: Collapse Visualization
%----------------------------------------------------------------------
\begin{figure}[h]
\centering
\begin{tikzpicture}[
    state/.style={circle, draw=blue!60, fill=blue!10, minimum size=7mm, drop shadow},
    ghost/.style={circle, draw=gray!30, fill=gray!5, minimum size=6mm, dashed},
    path/.style={-{Stealth}, thick, blue!80!black},
    ghostpath/.style={-{Stealth}, thin, gray!40, dashed},
    cone/.style={draw=orange!60, fill=orange!10, opacity=0.3},
    >=Stealth
]

% Future possibilities (cone) - DRAW FIRST so it's behind everything
\draw[cone, rounded corners]
    (0.8,0) -- (5.5,3.5) -- (5.5,-3.5) -- (0.8,0);

\node[font=\footnotesize, orange!80!black] at (4.8,0.2) {Time Cone};

% Current state
\node[state] (now) at (0,0) {$S_i$};

% Ghost futures (not taken) - wider spacing
\node[ghost] (g1) at (4.0,2.5) {};
\node[ghost] (g2) at (4.0,0.9) {};
\node[ghost] (g3) at (4.0,-0.9) {};
\node[ghost] (g4) at (4.0,-2.5) {};

\draw[ghostpath] (now) -- (g1);
\draw[ghostpath] (now) -- (g2);
\draw[ghostpath] (now) -- (g3);
\draw[ghostpath] (now) -- (g4);

% The chosen path (collapse)
\node[state] (next) at (4.0,0) {$S_{i+1}$};
\draw[path] (now) -- (next);

% Labels - positioned to avoid overlaps
\node[font=\footnotesize, blue!80!black, above=3pt] at (next.north) {Chosen};
\node[font=\scriptsize, gray, right=3pt] at (g1.east) {Discarded};
\node[font=\scriptsize, gray, right=3pt] at (g4.east) {Discarded};

% Collapse annotation - moved left to avoid overlapping diagram
\draw[->, ultra thick, red!70!black] (1.2,4.0) -- (1.2,0.7)
    node[midway, right=5pt, font=\footnotesize, align=left]
    {Collapse:\\Scheduler\\picks one};

\end{tikzpicture}
\caption{Collapse is the moment the Time Cone of possibilities narrows to a single chosen worldline.
The Scheduler evaluates all legal futures and commits to one, discarding the rest.}
\label{fig:collapse}
\end{figure}

This is the fundamental irreversibility of computation. Once a rewrite fires, you cannot un-fire it.
You cannot go back and try the other branch. You are committed. The state $S_i$ became $S_{i+1}$, and
there is no rewinding the tape. Classical systems hide this irreversibility behind abstractions like
"undo" and "rollback," but these are illusions built on top of careful state snapshotting.
\COMPUTER{} makes it explicit: Collapse is one-way. The worldline only grows forward.

\begin{center}
\rule{0.5\linewidth}{0.5pt}
\end{center}

\section{Worldlines Are Debugging}
\label{6.6-worldlines-are-debugging}

If execution is a path, then a bug is a wrong turn. Traditional debugging gives you
snapshots—stack traces, log lines, memory dumps. Static clues frozen in time. You get to see where
the program \emph{ended up}, but not how it got there. The path is invisible. You have to
reconstruct it from fragments, guessing at the sequence of events that led to the crash.

WARP-based debugging gives you the \emph{path} itself. You can examine where the worldline bent,
whether the rewrite was legal, and what structural condition misdirected it. You can replay the
execution tick by tick, inspecting every intermediate state. You can compare the buggy worldline to a
known-good worldline and see exactly where they diverge. You can measure Rulial Distance between them
and discover the precise coordinate where the universe took a bad turn.

Because worldlines and states are immutable, we can do something unheard of in classical systems:
\textbf{diff worldlines}. Not in the textual sense—this is not git diff on source code. This is a
geometric diff, a structural comparison of two paths through Rulial Space. We overlay the graphs,
cancel out identical nodes and edges, and highlight the differences. What remains is a map of
divergence: the exact subgraphs where the universes disagree.

\paragraph{The git bisect analogy.}
If a worldline is a commit history in the universe, then debugging it is a \emph{binary search on
reality}. You know the first tick was good. You know tick $N$ is bad. Binary search the worldline:
check tick $N/2$. If it's good, the bug is in the second half. If it's bad, the bug is in the first
half. Recurse until you isolate the single rewrite that introduced the error. Classical systems
cannot do this because they overwrite history. \COMPUTER{} cannot avoid doing this because
history \emph{is} the substrate. Every tick is recorded. Every state is preserved. The past is
immutable, and that makes it searchable.

Debugging becomes archaeology. You are not guessing—you are excavating. The evidence is all there,
buried in the worldline. You just have to dig in the right place.

\paragraph{The Rulial Heat Map.}
Diffing worldlines is not textual—it is geometric. Imagine overlaying the \emph{Happy Path} universe
and the \emph{Buggy Path} universe atop one another. Nodes and edges that match perfectly cancel
out; they become transparent. But any structural difference—a missing edge, an extra allocation, a
mutated subgraph—begins to glow.

\begin{quote}
Divergence has a temperature. Stable regions stay cool. Error regions burn hot.
\end{quote}

A worldline diff produces a \textbf{Rulial Heat Map}. Identical structure is transparent. Locally
altered structure glows warm orange. Cascading divergence burns deep red. Catastrophic
bifurcations—points where the two worldlines split into completely incompatible futures—appear as
white-hot fissures in possibility space.

Debugging stops being the reading of logs and becomes the reading of geometry. You see where the
universe itself is "hot" with error. A race condition is not a string of text in a log file. It is a
thermal scar on the shape of execution. It is a visible wound in the geometry, and you can point to
it and say: there. That is where the worldline broke.

This is visceral, visual, unforgettable. Engineers will feel this. They will see the heat map and
understand, in their gut, what went wrong. This is not abstract theory. This is debugging with your
eyes open.

%----------------------------------------------------------------------
% TikZ: Rulial Heat Map
%----------------------------------------------------------------------
\begin{figure}[h]
\centering
\begin{tikzpicture}[
    node/.style={circle, draw=black, minimum size=7mm, inner sep=1pt, drop shadow},
    same/.style={circle, draw=black!30, fill=black!10, minimum size=7mm},
    warm/.style={circle, draw=orange!80!black, fill=orange!40, minimum size=7mm, drop shadow},
    hot/.style={circle, draw=red!80!black, fill=red!50, minimum size=7mm, drop shadow},
    whitehot/.style={circle, draw=red!90!black, fill=white, line width=1pt, minimum size=7mm, drop shadow},
    edge/.style={-{Stealth}, thick},
    ghost/.style={draw=none, fill=none}
]

% "Correct" baseline graph (left)
\node[same] (c1) at (-3,0.5) {\tiny A};
\node[same] (c2) at (-3,-0.5) {\tiny B};
\node[same] (c3) at (-4,0) {\tiny C};
\node[same] (c4) at (-2,0) {\tiny D};

\draw[edge, gray!50] (c1) -- (c3);
\draw[edge, gray!50] (c3) -- (c2);
\draw[edge, gray!50] (c1) -- (c4);
\draw[edge, gray!50] (c2) -- (c4);

\node[ghost] at (-3,-1.4) {Correct Worldline Slice};

% "Buggy" worldline overlay (right)
\node[same] (b1) at (3,0.5) {\tiny A};
\node[warm] (b2) at (3,-0.5) {\tiny B};
\node[hot] (b3) at (2,0) {\tiny C};
\node[whitehot] (b4) at (4,0) {\tiny D};

% Edges
\draw[edge, gray!50] (b1) -- (b3);
\draw[edge, orange!80!black, thick] (b3) -- (b2);
\draw[edge, red!80!black, thick] (b1) -- (b4);
\draw[edge, red!90!black, ultra thick] (b2) -- (b4);

\node[ghost] at (3,-1.4) {Buggy Worldline Slice};

% Heat legend - increased spacing from 8pt to 15pt
\node[same] at (0,-2.0) {};
\node[ghost, right=15pt of {(0,-2.0)}] {\footnotesize Identical (transparent)};

\node[warm] at (0,-2.6) {};
\node[ghost, right=15pt of {(0,-2.6)}] {\footnotesize Local divergence (warm)};

\node[hot] at (0,-3.2) {};
\node[ghost, right=15pt of {(0,-3.2)}] {\footnotesize Strong divergence (hot)};

\node[whitehot] at (0,-3.8) {};
\node[ghost, right=15pt of {(0,-3.8)}] {\footnotesize Catastrophic divergence (white-hot)};

\end{tikzpicture}
\caption{Rulial Heat Map: comparing the happy-path universe to the buggy universe.
Identical structure cancels out; divergent structure radiates heat.}
\label{fig:rulial-heat-map}
\end{figure}

%----------------------------------------------------------------------
% TikZ: Multi-Frame Rulial Heat Animation
%----------------------------------------------------------------------
\begin{figure}[h]
\centering
\begin{tikzpicture}[
    same/.style={circle, draw=black!30, fill=white, minimum size=7mm},
    warm/.style={circle, draw=orange!80!black, fill=orange!40, minimum size=7mm, drop shadow},
    hot/.style={circle, draw=red!80!black, fill=red!50, minimum size=7mm, drop shadow},
    whitehot/.style={circle, draw=red!90!black, fill=white, line width=1pt, minimum size=7mm, drop shadow},
    edge/.style={-{Stealth}, thin, gray!60},
    frame/.style={rectangle, rounded corners, draw=gray!50, fill=white, inner sep=12pt, drop shadow}
]

% Frame 1: Early divergence - scaled up
\node[frame] (f1) at (-5,0) {
    \begin{tikzpicture}[scale=1.0]
        \node[same] (a1) at (0,0.6) {};
        \node[same] (a2) at (-0.6,0) {};
        \node[same] (a3) at (0.6,0) {};
        \node[warm] (a4) at (0,-0.6) {};
        \draw[edge] (a1) -- (a2);
        \draw[edge] (a1) -- (a3);
        \draw[edge] (a2) -- (a4);
        \draw[edge] (a3) -- (a4);
    \end{tikzpicture}
};

\node[font=\scriptsize, above] at (f1.north) {$t_0$: Local anomaly};

% Frame 2: Spreading heat - scaled up
\node[frame] (f2) at (0,0) {
    \begin{tikzpicture}[scale=1.0]
        \node[warm] (b1) at (0,0.8) {};
        \node[warm] (b2) at (-0.7,0.2) {};
        \node[hot]  (b3) at (0.7,0.2) {};
        \node[hot]  (b4) at (-0.4,-0.6) {};
        \node[same] (b5) at (0.6,-0.7) {};
        \draw[edge] (b1) -- (b2);
        \draw[edge] (b1) -- (b3);
        \draw[edge] (b2) -- (b4);
        \draw[edge] (b3) -- (b4);
        \draw[edge] (b4) -- (b5);
    \end{tikzpicture}
};

\node[font=\scriptsize, above] at (f2.north) {$t_1$: Heat propagates};

% Frame 3: Catastrophic divergence - scaled up
\node[frame] (f3) at (5,0) {
    \begin{tikzpicture}[scale=1.0]
        \node[hot]     (c1) at (0,1.0) {};
        \node[whitehot](c2) at (-0.7,0.2) {};
        \node[whitehot](c3) at (0.7,0.2) {};
        \node[hot]     (c4) at (-0.4,-0.6) {};
        \node[hot]     (c5) at (0.6,-0.7) {};
        \draw[edge] (c1) -- (c2);
        \draw[edge] (c1) -- (c3);
        \draw[edge, red!80!black, thick] (c2) -- (c4);
        \draw[edge, red!80!black, thick] (c3) -- (c4);
        \draw[edge, red!90!black, ultra thick] (c4) -- (c5);
    \end{tikzpicture}
};

\node[font=\scriptsize, above] at (f3.north) {$t_2$: Full-blown failure region};

% arrows
\draw[->, thick, gray!70] (f1.east) -- (f2.west)
    node[midway, above=1pt, font=\scriptsize] {propagation};
\draw[->, thick, gray!70] (f2.east) -- (f3.west)
    node[midway, above=1pt, font=\scriptsize] {runaway};

\end{tikzpicture}
\caption{Temporal evolution of Rulial Heat. A small local deviation at $t_0$
becomes a structurally significant divergence at $t_1$, and by $t_2$ the bug
has carved out an entire hot region in the universe. Debugging is the act of
locating and cooling this region.}
\label{fig:rulial-heat-animation}
\end{figure}

\begin{mdframed}[linewidth=1pt, linecolor=black!50, innertopmargin=8pt, innerbottommargin=8pt]
\noindent\textbf{FOR THE NERDS™: Rulial Heat as Divergence}

\medskip

Let $(S_i^{\mathrm{good}})_{i=0}^N$ and $(S_i^{\mathrm{bug}})_{i=0}^N$ be two worldlines
with the same initial state $S_0$ and identical tick index set.

For each tick $i$, define the \textbf{instantaneous structural divergence}:

\[
\Delta_i
= S_i^{\mathrm{good}} \triangle S_i^{\mathrm{bug}},
\]

where $\triangle$ denotes the symmetric difference over the set of
\emph{typed graph elements} (nodes, edges, and their adjacencies).

We define the \textbf{Rulial Heat} at tick $i$ as:

\[
H_i = |\Delta_i|.
\]

Intuitively, when $H_i = 0$, the universes are structurally identical at tick $i$. Larger $H_i$ values indicate more structural disagreement, corresponding to more "heat" in the system.

We can then define:

\[
H_{\max} = \max_{0 \le i \le N} H_i,
\quad
H_{\mathrm{tot}} = \sum_{i=0}^N H_i,
\]

measuring, respectively, the \textbf{peak heat} and the \textbf{integrated heat} of the divergence.

\begin{quote}
Rulial Heat is how much the buggy universe disagrees with the correct one, measured
not in log lines but in altered structure.
\end{quote}

\end{mdframed}

\paragraph{The Observer Effect Is Dead.}
In classical debugging, attaching a debugger \emph{changes} the program. Adding
a \texttt{printf()} shifts timing, scares away race conditions, and collapses
the very bug you're trying to study. These "Heisenbugs" survive by reacting to
observation. They are like quantum particles that refuse to be measured without changing their state.

In \COMPUTER{}, the Observer Effect is dead.

Because the Scheduler ignores the wall clock entirely, you can pause the
universe for a thousand years between ticks, inspect every atom, and when you resume
execution the next rewrite will be \emph{identical}. Timing cannot influence causality. The Scheduler
does not care how long you spent staring at the graph. It does not care if you attached a debugger.
It does not care if you printed the entire state to a log file. When you unpause, it will pick the
same match, apply the same rewrite, and produce the same successor state. Every. Single. Time.

\begin{quote}
The Heisenbug cannot run. It cannot hide. Time has no power here.
\end{quote}

Determinism makes debugging safe again: observing the universe no longer disturbs the universe. The
worldline is what it is. You can study it, diff it, replay it, dissect it—and it will not flinch.

%----------------------------------------------------------------------
% TikZ: Heisenbug Extinction
%----------------------------------------------------------------------
\begin{figure}[h]
\centering
\begin{tikzpicture}[
    world/.style={-{Stealth}, thick},
    altworld/.style={-{Stealth}, thick, dashed},
    tick/.style={circle, fill=black, inner sep=1pt},
    labelstyle/.style={font=\scriptsize},
    >=Stealth
]

% Classical system (top)
\node[labelstyle] at (-5,2.4) {Classical System};

\draw[world, gray!70, line width=0.8pt] (-5,2.0) -- (-0.5,2.0);
\draw[altworld, red!80!black, line width=0.8pt] (-5,1.7) -- (-2.5,1.7) .. controls (-1.5,1.7) and (-0.8,1.2) .. (0,0.8);

\node[tick, minimum size=2.5pt] (c0) at (-5,2.0) {};
\node[tick, minimum size=2.5pt] (c1) at (-3.5,2.0) {};
\node[tick, minimum size=2.5pt] (c2) at (-2,2.0) {};
\node[tick, minimum size=2.5pt] (c3) at (-0.5,2.0) {};

\node[labelstyle, above] at (-3.5,2.0) {$t_1$};
\node[labelstyle, above] at (-2,2.0) {$t_2$};

% breakpoint marker
\draw[red!80!black, thick] (-2,2.25) -- (-2,1.75);
\node[labelstyle, above, red!80!black] at (-2,2.35) {Debugger attached};

\node[labelstyle, align=left] at (2.5,1.8)
    {\scriptsize Same input, different timing\\\scriptsize $\Rightarrow$ different outcome (Heisenbug)};

% CΩMPUTER system (bottom)
\node[labelstyle] at (-5,-0.4) {\COMPUTER{}};

\draw[world, blue!80!black, line width=0.8pt] (-5,-0.8) -- (-0.5,-0.8);
\draw[altworld, blue!40!black, densely dotted, line width=0.8pt] (-5,-1.1) -- (-0.5,-1.1);

\node[tick, minimum size=2.5pt] (d0) at (-5,-0.8) {};
\node[tick, minimum size=2.5pt] (d1) at (-3.5,-0.8) {};
\node[tick, minimum size=2.5pt] (d2) at (-2,-0.8) {};
\node[tick, minimum size=2.5pt] (d3) at (-0.5,-0.8) {};

\node[labelstyle, above] at (-3.5,-0.8) {$t_1$};
\node[labelstyle, above] at (-2,-0.8) {$t_2$};

% breakpoint marker
\draw[gray!70, thick] (-2,-0.55) -- (-2,-1.05);
\node[labelstyle, above, gray!70] at (-2,-0.45) {Paused arbitrarily};

\node[labelstyle, align=left] at (2.5,-0.8)
    {\scriptsize Wall-clock ignored, Scheduler canonical\\\scriptsize $\Rightarrow$ same outcome, no Heisenbug};

\end{tikzpicture}
\caption{In classical systems, attaching a debugger perturbs timing and can change the outcome
(a Heisenbug). In \COMPUTER{}, the Scheduler ignores time, so observation cannot alter the
worldline. The Observer Effect is dead.}
\label{fig:heisenbug-extinction}
\end{figure}

\begin{center}
\rule{0.5\linewidth}{0.5pt}
\end{center}

\section{Worldlines Are Optimization}
\label{6.7-worldlines-are-optimization}

Optimization is nothing but \textbf{geodesic alignment}. This is not a metaphor. When you optimize
code, you are performing a geometric operation: you are unbending the worldline, reducing its
curvature, pushing it closer to the shortest path through Rulial Space. Every classical optimization
trick—constant folding, dead code elimination, loop unrolling, inlining—is an act of straightening.

Consider dead code elimination. When the compiler removes a function that is never called, it is not
just "cleaning up" the source. It is removing an entire detour from the worldline. That function
represented potential rewrites—nodes that could have been allocated, edges that could have been
traversed. By proving the function is unreachable, the compiler deletes those possibilities entirely.
The worldline no longer has to bend through that region of Rulial Space. It gets shorter, straighter,
faster.

Inlining works by collapsing wormholes. A function call introduces overhead: push arguments onto the
stack, jump to a new address, execute the body, pop the return value, jump back. Each of these steps
is a rewrite. Each rewrite bends the worldline. Inlining replaces the call with the body, eliminating
the jump. The worldline now takes a direct route through the logic, skipping the detour entirely. The
Rulial Distance between input and output shrinks.

Memoization creates shortcuts. If a function has already computed \texttt{f(x)}, and you call
\texttt{f(x)} again, memoization returns the cached result instead of recomputing. This is not just
"saving time"—it is creating a wormhole in Rulial Space. The second call avoids retracing the entire
path through the computation. Instead, it takes a single rewrite: "look up the answer." The worldline
bends around the expensive region, exploiting the fact that the geometry allows this shortcut.

Fast code feels elegant because it is geometrically direct. When you read well-optimized code, you
sense its economy of motion. There is no waste. Every line has a purpose. Every function does exactly
what it must, and nothing more. This is not accident. It is geometry. The worldline is close to the
geodesic, and you can feel it in the structure of the code.

Conversely, slow code feels wrong because it is wandering through high-curvature regions for no good
reason. It allocates and frees. It checks the same condition twice. It recomputes values it already
knows. Each of these inefficiencies is a visible bend in the worldline, and experienced engineers can
spot them instantly. Not because they memorized a list of anti-patterns, but because they have
internalized the geometry. They can see when the path is curved, and they know how to straighten it.

\subsection{The Profiler as Curvature Detector}

A profiler does not "find slow code." It finds regions of high curvature. When a profiler highlights
a function that consumes 80\% of runtime, it is telling you: this is where the worldline is bending
the most. This is where the path is farthest from the geodesic. If you want to make the program
faster, straighten this part of the worldline first.

Optimization, then, is not about random tweaks. It is a deliberate, geometric process. You measure
the worldline's curvature. You identify the steepest bends. You apply transformations—refactors,
algorithmic improvements, caching strategies—that reduce those bends. You re-measure. You iterate.
Each cycle brings the worldline closer to the geodesic. Each cycle makes the program faster, not
because you tried harder, but because you aligned the execution with the natural geometry of the
problem.

\begin{mdframed}[linewidth=1pt, linecolor=black!50, innertopmargin=8pt, innerbottommargin=8pt]
\noindent\textbf{FOR THE NERDS™: Optimization as Discrete Gradient Descent}

\medskip

Fix an input state $S_{\mathrm{in}}$ and a desired output state $S_{\mathrm{out}}$.
Consider the set $\mathcal{W}$ of all legal worldlines:

\[
\mathcal{W} =
\left\{
(S_i)_{i=0}^N \ \middle|\
S_0 = S_{\mathrm{in}},\ S_N = S_{\mathrm{out}},\ S_{i+1} = \mathrm{step}(S_i)
\right\}.
\]

Define a \textbf{cost functional} over worldlines:

\[
\mathcal{C}((S_i)_{i=0}^N)
= \sum_{i=0}^{N-1} c(S_i, S_{i+1}),
\]

where $c$ is a local cost. We can choose $c = 1$ to measure path length (equivalent to Rulial Distance), or $c = \mathcal{F}(S_i)$ to incorporate friction-weighted cost, or we can use a weighted combination that accounts for resource usage such as memory consumption, I/O operations, and other system metrics.

Optimization is then:

\[
\text{Find } (S_i^\star) \in \mathcal{W} \text{ minimizing } \mathcal{C}.
\]

Compiler transforms and refactors can be viewed as \textbf{discrete gradient steps}
on $\mathcal{C}$: local rewrites to the rule set or program graph that strictly
decrease the expected cost along the induced worldline.

\begin{quote}
You are not "making the code faster." You are performing gradient descent
on the manifold of possible universes, aligning your worldline with the geodesic.
\end{quote}

\end{mdframed}

\begin{center}
\rule{0.5\linewidth}{0.5pt}
\end{center}

\begin{mdframed}[linewidth=1pt, linecolor=black!50, innertopmargin=8pt, innerbottommargin=8pt]
\noindent\textbf{FOR THE NERDS™: Worldlines and Lambda Calculus Reduction}

\medskip

If you squint, a worldline is a reduction sequence in Lambda Calculus. Each rewrite is a kind of
redex collapse. The Scheduler picks which redex to reduce next, just as a reduction strategy picks
leftmost-outermost or call-by-value. The analogy is seductive, and it holds up to a point.

But it breaks quickly. WARP graph reductions operate in a metric space with measurable distances
between states. Lambda calculus has no notion of "how far" one term is from another. WARP graphs
support explicit structure persistence: nodes and edges have stable identities that survive rewrites.
Lambda terms are ephemeral; alpha-equivalence erases identity. WARP graphs allow recursive edges,
cycles, and multi-scale nested subgraphs. Lambda calculus is strictly tree-structured, with no native
support for sharing or cycles.

The key point is this: the \textbf{Scheduler is the reduction strategy}. Not metaphorically—literally.
Just as lambda calculi choose between call-by-name, call-by-value, or call-by-need, the Scheduler
chooses which match to collapse next. The difference is that WARP space is curved, structured, and
globally coherent. Lambda calculus is flat and local. WARP worldlines have geometry. Lambda reductions
do not.

This is why we needed \COMPUTER{} in the first place. Lambda calculus is a beautiful model of
pure computation, but it has no concept of distance, friction, or curvature. It cannot explain why
some programs are fast and others are slow. It cannot model technical debt or refactoring difficulty.
It cannot visualize the path through state space. \COMPUTER{} can.

\end{mdframed}

\begin{center}
\rule{0.5\linewidth}{0.5pt}
\end{center}

\section{Transition: From Worldlines to Neighborhoods}
\label{6.8-transition}

We now understand individual execution paths. We know what a worldline is, why it is deterministic,
how it bends, and what it means for it to be straight. But systems are not just lines—they are
\emph{fields}. A single worldline is one trajectory through a vast space of possibilities. To
understand real computation—to answer questions about robustness, flexibility, and refactorability—we
must study the \textbf{space around the worldline}.

Which nearby futures are viable? If you perturb the input slightly, does the worldline stay close, or
does it veer off into disaster? Which states are "adjacent" in the sense that a small refactor can
reach them, and which require massive structural surgery? Why do some systems feel flexible—easy to
modify, easy to extend—while others feel rigid and brittle, where every change threatens to shatter
the whole edifice?

These questions pull us into the domain of local geometry—\textbf{Neighborhoods}—the topic of
Chapter 7. Where worldlines are paths, neighborhoods are regions. Where worldlines ask "what did the
system do?", neighborhoods ask "what else could it have done?" We leave the certainty of the single
chosen path and step into the fog of nearby possibilities, armed with the geometry we have built so
far.

%----------------------------------------------------------------------
% Deterministic Scheduler Specification
%----------------------------------------------------------------------
\clearpage

\begin{center}
\rule{0.7\linewidth}{1.5pt}
\vspace{0.3cm}

{\Large \textbf{APPENDIX: FORMAL SPECIFICATIONS}}

\vspace{0.2cm}
\rule{0.7\linewidth}{1.5pt}
\end{center}

\vspace{0.5cm}

\noindent\textit{The following sections provide the formal mathematical machinery and implementation contracts
for those building C$\Omega$MPUTER runtimes. These specifications are not required to understand the
conceptual content of Chapter 6, but they define the exact semantics that implementations must preserve.}

\vspace{0.5cm}

\section{Deterministic Scheduler Specification}
\label{sec:scheduler-spec}

In this section we make precise what was stated informally in
Chapter~\ref{chapter-6-worldlines-execution-as-geodesics}: that the
Scheduler induces a \emph{canonical} next-step transition for every
reachable universe state.

\section{Axioms and Notation}

Let:

\begin{itemize}
    \item $\mathcal{S}$ be the set of all valid WARP graph universe states.
    \item $\mathcal{R}$ be a finite set of DPO rewrite rules.
    \item Each rule $r \in \mathcal{R}$ has a unique identifier $\mathrm{id}(r)$
          drawn from a totally ordered set $(\mathbb{R}_{\text{id}}, <_{\text{rule}})$.
    \item Each rule $r$ is given by a production $(L_r \xleftarrow{l_r} K_r \xrightarrow{r_r} R_r)$
          in the usual DPO sense.
\end{itemize}

For a given state $S \in \mathcal{S}$, define the set of \textbf{legal matches}:

\[
\mathsf{Matches}(S)
= \left\{ (r, \mu) \,\middle|\,
    r \in \mathcal{R},\;
    \mu : L_r \hookrightarrow S\ \text{is a monomorphism satisfying all gluing conditions}
    \right\}.
\]

We assume:

\begin{enumerate}
    \item Each universe state $S$ has a stable node and edge identity scheme,
          inducing a total order $<_{\text{id}}$ over all node and edge IDs.
    \item For each match $(r, \mu)$ we can extract a finite, ordered tuple of IDs
          touched by the match:
          \[
          \mathsf{ids}(r, \mu) = (i_1, i_2, \dots, i_k)
          \]
          where $(i_1 <_{\text{id}} i_2 <_{\text{id}} \dots <_{\text{id}} i_k)$.
    \item We assume the existence of a \emph{canonical embedding hash}
          \[
          h_S : \mathsf{Matches}(S) \to \mathbb{H}
          \]
          into a totally ordered set $(\mathbb{H}, <_{\text{hash}})$ such that
          distinct matches induce distinct hashes. (In practice this can be a
          deterministic hash of $(\mathrm{id}(r), \mathsf{ids}(r,\mu), S)$.)
\end{enumerate}

\section{Ordering Matches}

We now define a strict total order $<_{\text{match}}$ on $\mathsf{Matches}(S)$
which the Scheduler will use for selection.

For two matches $m_1=(r_1,\mu_1)$ and $m_2=(r_2,\mu_2)$ in $\mathsf{Matches}(S)$,
we say $m_1 <_{\text{match}} m_2$ iff the following lexicographic chain holds:

\begin{enumerate}
    \item \textbf{Rule priority:}
          \[
          \mathrm{id}(r_1) <_{\text{rule}} \mathrm{id}(r_2),
          \]
          or, if equal,
    \item \textbf{Structural footprint:}
          \[
          \mathsf{ids}(r_1,\mu_1) <_{\text{lex}} \mathsf{ids}(r_2,\mu_2),
          \]
          or, if equal,
    \item \textbf{Canonical embedding hash:}
          \[
          h_S(m_1) <_{\text{hash}} h_S(m_2).
          \]
\end{enumerate}

Here $<_{\text{lex}}$ is the usual lexicographic order induced by $<_{\text{id}}$.

\begin{quote}
The third component ($h_S$) only matters in pathological cases where two matches
have identical rule IDs and identical ordered ID-tuples. It serves as a
tie-breaker of last resort, ensuring a strict total order on matches.
\end{quote}

\section{Conflict Elimination}

Two matches $m_1 = (r_1, \mu_1)$ and $m_2 = (r_2, \mu_2)$ are said to
\textbf{conflict} if their images overlap:

\[
\mathrm{im}(\mu_1) \cap \mathrm{im}(\mu_2) \neq \varnothing.
\]

Given $\mathsf{Matches}(S)$, define the \emph{conflict-free subset} selected
by leftmost-priority elimination:

\[
\mathsf{cf}(S) = \text{the unique subset obtained by:}
\]

\begin{enumerate}
    \item Sorting $\mathsf{Matches}(S)$ in ascending $<_{\text{match}}$ order.
    \item Scanning from smallest to largest, greedily including a match if and only if
          it does not conflict with any match already included.
\end{enumerate}

\medskip

In the worldline semantics used in Chapter~\ref{chapter-6-worldlines-execution-as-geodesics}
we apply exactly \emph{one} rewrite per tick, so we do not exploit the full parallelism of
$\mathsf{cf}(S)$. However, the same machinery can later be reused to define
\emph{parallel} steps.

\section{The Scheduler Function}

We can now define the Scheduler as a partial function
\[
\sigma : \mathcal{S} \rightharpoonup \mathsf{Matches}(S)
\]
given by:

\[
\sigma(S) =
\begin{cases}
\text{the $<_{\text{match}}$-minimal element of } \mathsf{cf}(S), & \text{if } \mathsf{cf}(S) \neq \varnothing, \\
\bot, & \text{otherwise,}
\end{cases}
\]
where $\bot$ denotes that no further rewrites are possible (a halting state).

\section{Determinism Guarantee}

Given an initial state $S_0$, the worldline $(S_i)_{i \ge 0}$ defined by iterative application of
$\mathrm{step}(S)$ is uniquely determined. This guarantees that different machines cannot produce
different universes.

\section{Runtime Contract for Deterministic Execution}
\label{appendix:runtime-contract}

This appendix specifies the \emph{minimal invariants} required of any
implementation of the \COMPUTER{} runtime. These are not guidelines.
They are hard constraints. Violate them and you do not have a deterministic
universe anymore—you have a slot machine wearing a lab coat.

This contract exists to ensure that the worldline semantics of
Chapter~\ref{chapter-6-worldlines-execution-as-geodesics} hold exactly,
byte-for-byte, across replays, reboots, architectures, and compilers.

\section{Stable Identity Invariant}

Every node and edge in a WARP graph must carry a \textbf{stable, opaque identity}
drawn from a totally ordered set. These IDs \emph{must not} encode memory addresses, timestamps,
allocation order, or any machine-dependent data. Furthermore, IDs \emph{must survive} serialization,
deserialization, snapshotting, garbage collection, graph compaction, and transport without corruption
or reordering. The order $<_{\text{id}}$ must be global and total across all system components.

\begin{quote}
If IDs drift, reorder, or depend on runtime artifacts, the universe fractures.
\end{quote}

\section{Deterministic Match Enumeration}

Given a state $S$, every legal match $(r,\mu)$ must be constructed in a \textbf{deterministic}
enumeration order, expressed as a tuple of IDs \emph{sorted by} $<_{\text{id}}$, and completely
free of any dependence on hash randomization, memory layout, or concurrency artifacts.

\begin{quote}
If the set of matches is nondeterministic, the Scheduler becomes a liar.
\end{quote}

\section{Lexicographic Rule Ordering Must Be Canonical}

Each rewrite rule $r$ must have a globally unique rule ID $\mathrm{id}(r)$ and a total ordering on
IDs. This means rule identities must not depend on file paths, loader order, timestamps, random
seeds, or compiler heuristics. Rule identities must not be generated at runtime; they \emph{must}
be fixed at compile time or definition time to ensure reproducibility across all execution contexts.

\section{Conflict Resolution Must Be Purely Structural}

When eliminating overlapping matches, overlap must be defined purely as intersection of
$\mathrm{im}(\mu)$ sets, and the greedy leftmost-priority rule must be used exactly as defined in
Section~\ref{sec:scheduler-spec}. No heuristic such as "prefer larger matches" or "prefer denser
subgraphs" may be substituted under any circumstances.

\begin{quote}
If conflict-resolution is heuristic, replay is impossible.
\end{quote}

\section{Canonical Embedding Hash Must Be Deterministic}

The embedding hash $h_S(m)$ must satisfy strict determinism requirements. For a given $(S,m)$, the
hash must be \textbf{bit-identical} across all runs, and hash functions must not use randomized
salts. Hash inputs must derive only from three sources: first, $\mathrm{id}(r)$; second,
$\mathsf{ids}(r,\mu)$; and third, the structural context of $S$ reachable from the match. The
implementation must define explicit normalization rules for serialization to ensure hash stability.

\begin{quote}
If the hash isn't stable, your Scheduler is choosing futures with a coin flip.
\end{quote}

\section{DPO Construction Must Be Canonical}

Given $(r,\mu)$, node and edge deletions must occur in an order independent of container iteration
details, and newly created IDs must be drawn in deterministic order. The embedding of $R$ into $S'$
must use canonical port ordering, and graph rewriting must not depend on concurrent data structures
whose iteration order varies across platforms or runs.

\begin{quote}
Different machines cannot produce different universes.
\end{quote}

\section{Time Must Not Exist}

No rewrite \emph{may} depend on wall-clock time, monotonic clocks, random number generators (unless
seeded deterministically), thread scheduling order, input arrival timing, or system clock drift. Time
must be eliminated as a factor in rewrite selection.

\begin{quote}
\COMPUTER{} does not know what day it is.
\end{quote}

\section{8. No Hidden State Allowed}

All state influencing rewrite decisions must be either stored explicitly in the WARP graph or encoded
in the rule ID and match structure. No implementation-defined caches, memo-tables, or thread-local
scratchpads may influence rule applicability.

\begin{quote}
If it isn't in the graph, it isn't in the universe.
\end{quote}

\section{9. Serialization Must Be Lossless and Deterministic}

Any serialization round-trip $S \to \text{bytes} \to S'$ must satisfy:

\[
S' = S.
\]

This includes graph structure, node and edge identities, port ordering, rule metadata, and provenance
annotations. No aspect of the system state may be lost or corrupted during serialization and
deserialization.

\section{10. Replay Must Be Perfect}

Given a worldline $(S_i)_{i=0}^N$:

\[
S'_0 = S_0 \quad \Rightarrow \quad S'_i = S_i \quad \forall i.
\]

If this property does not hold, then one of four failures has occurred: either the Scheduler is
nondeterministic, or the graph is not canonical, or serialization is flawed, or the implementation
has violated this contract in some other way.

\section{11. Fast Is Optional. Determinism Is Not.}

A runtime may parallelize match enumeration, cache embeddings, fuse rule applications, compact graphs,
and JIT-compile rule kernels. However, all such optimizations must preserve the exact
\textbf{Scheduler-visible surface}:

\[
\sigma(S) \text{ must not change.}
\]

If it does, the optimization is illegal. Parallelism is an implementation detail; Determinism is the law.

\begin{center}
\rule{0.7\linewidth}{0.7pt}
\end{center}

\section*{The One-Sentence Contract}

\begin{quote}
\Large\textbf{
Everything in the runtime may change—except the order in which the universe chooses its next legal rewrite.
}
\end{quote}

Break this, and you break determinism. Preserve it, and you inherit all the geometry.
